\section{Genauigkeit der Verfahren}
\label{sec:genauigkeit}

Die Laufversuche liefern eine Bruchlinie nach dem Prüfstandskriterium, der Nachweis nach Stommel et al. eine Lastwechselzahl bis zum ersten makroskopischen Anriss. Nach \secref{sec:auswertung_laufversuche} gilt für jeden Lauf $N_{\mathrm{A}} \le N_{\mathrm{K}} \le N_{\mathrm{B}}$, so dass der Vergleich in drei Stufen geführt wird. In der ersten Stufe steht die Vorhersage gegen die Lastwechselzahl am Ende des Laufs. Mit der Stufe W2 liegt sie bei $8$, $12$ und $16$\,Nm um die Faktoren $75$, $270$ und $820$ unter der gemessenen Linie für $50$\,\%, mit W1 und mit allen Hill-Varianten weiter darunter. Die Vorhersage liegt damit an jeder Stufe auf der sicheren Seite, und der Abstand wächst mit dem Moment. Dieses Verhältnis ist die Gesamtreserve aus der Rissfortschrittsphase und der Reserve der Methode; beide lassen sich ohne Anrissdaten nicht trennen. In der zweiten Stufe stehen die Anrissfenster aus \secref{sec:woehlerlinien}: Bei $8$\,Nm wurden die ersten Anrisse zwischen $9$ und $31$\,\% der Laufzeit sichtbar, so dass der Abstand der Methode zu ihrer eigenen Versagensdefinition dort auf den Faktor $9$ bis $22$ schrumpft. Bei $12$ und $16$\,Nm liegt nur die erste Kontrolle mit Befund vor, zwischen $55$ und $93$\,\% der Laufzeit; der Anriss kann dort früher liegen, und der Abstand ist nach unten offen. Dass bei $8$\,Nm mehr als zwei Drittel der Laufzeit mit sichtbaren Rissen vergingen, zeigt, wie weit die Rissfortschrittsphase bei diesem Werkstoff die Lebensdauer bis zum Bruch trägt. Ein Nachweis, der beim Anriss endet, kann die Bruchlinie deshalb nicht treffen; wer bis zum Bruch auslegen will, braucht entweder eine Rissfortschrittsreserve oder eine am Versuch kalibrierte Linie.

Die dritte Stufe ist die lokale Wöhlerlinie am kritischen Knoten. Trägt man die gemessenen Lastwechselzahlen über der FE-Spannung an C auf, so folgt aus den Werten der \tabref{tab:verfahrensvergleich} eine doppeltlogarithmische Steigung von $k = 12$ mit den Spannungen der Stufe W2 und von $10$ mit W4. Mit W1 und W3 sind es $7{,}6$ und $7{,}4$, weil die Spannung der plastischen Stufen unterproportional zum Moment wächst. Die synthetische Wöhlerlinie des Nachweises trägt bei PA46-GF30 die Neigung $14{,}8$ bis $15{,}8$ aus den Klassentabellen und die Wechselfestigkeit $20$\,MPa bei $10^{7}$ Lastwechseln. Die Steigung der plastischen Stufen liegt damit nahe an der des Buches, die Lage der Linie nicht: Der Zahn erträgt bei $8$\,Nm mit $38$\,MPa größter Zugspannung an der Oberfläche $2{,}7 \cdot 10^{7}$ Lastwechsel bis zum Bruch. Eine am Versuch kalibrierte Linie würde deshalb vor allem die Lage anheben und die Neigung behalten; sie bleibt eine Bruchlinie und ist als eigene Erweiterung außerhalb der Buchdefinition zu benennen. \Note{PH: lokale Wöhlerlinie mit Amplitude und Mittelspannung des kritischen Knotens je Stufe (E8), Zahlen nach dem Lauf des Skripts.}

Mehrere Eigenschaften der Methode tragen zu ihrer Reserve bei. Die Isotropievorschrift nimmt die Kurve in Faserrichtung als isotrope Referenz, während die Klassenfaktoren der verstärkten Werkstoffe die Querschwäche pauschal abbilden. Der kritische Knoten liegt auf der Druckseite der Fußrundung mit $R \approx -0{,}4$; dort rechnet die Methode über den Überlastfall F2 und den Mittelspannungsfaktor, obwohl ihre Festigkeitswerte für schwellende Beanspruchung aufgenommen wurden. Die Spannungen sind Größtwerte an der Oberfläche nahe der Stirnseite des Zahns, bei $z = -6{,}0$\,mm, wo die freie Stirnfläche die Fußrundung schneidet. Die Werkstoffkarte gilt bei $80$\,\textdegree C, der Zahn lief bei $16$\,Nm nach \secref{sec:temperatur} mit $84$ bis $88$\,\textdegree C wärmer; das senkt die Festigkeit und erklärt einen Teil des mit dem Moment wachsenden Abstands nicht, sondern vergrößert ihn. Die Belastungsfrequenz von $50$\,Hz liegt über der quasistatischen Karte; die Dehnratenkarten erhöhen Steifigkeit und Festigkeit zugleich, ihre Wirkung auf die Vorhersage ist nach \secref{sec:stufen_raten_koerper} klein gegen den Sprung zwischen den Stufen. Gegen S-Life Plastics ist die Umsetzung an Spannungen, statischem Nachweis, Amplitude und Festigkeit bestätigt; der Faktor $2{,}5$ in der zyklischen Auslastung ist allein die Kappe der zulässigen Amplitude. Welche Lesart das Buch vorschreibt, ändert die Vorhersage um diesen Faktor, aber nicht die Größenordnung des Abstands zum Versuch.

Für die Nennspannung nach VDI 2736 Blatt 2, RIKOR und die FVA-Workbench steht der Vergleich in \tabref{tab:verfahrensvergleich} aus. \Note{PH: Einordnung von VDI 2736, RIKOR und Workbench gegen die FE-Spannung und die Bruchlinie, sobald die Werte vorliegen; Steigung der Nennspannung proportional zur Umfangskraft gegen $k = 4{,}40$ der Bruchlinie.} Die Verteilung der Lastwechselzahlen beschreiben die logarithmische Normalverteilung und die Weibull-Verteilung mit neun Ergebnissen gleich gut; die Faktoren von $50$ auf $10$\,\% liegen je nach Weg zwischen $0{,}60$ und $0{,}74$, nach VDI 2736 Blatt 4 am höchsten. Für die Einordnung des Werkstoffs trägt der Vergleich auf demselben Prüfstand: Bei $16$\,Nm liegt PA46-GF30 um den Faktor $1{,}6$ bis $2{,}1$ über POM und bei $14$ bis $18$\,Nm unter PEEK, dessen Räder im Jahr 2023 jedoch an der Flanke und nicht am Fuß versagten.

\section{Anwendbarkeit und Modellgrenzen}
\label{sec:modellgrenzen}

Die Methode nach Stommel et al. ist in der veröffentlichten Form anwendbar, wenn ihre Versagensdefinition und ihre Isotropievorschrift mitgenannt werden. Unter den vier Werkstoffstufen ist W2 die einzige, deren statischer Nachweis an allen Fußknoten besteht und deren Lebensdauer mit dem Moment in plausibler Ordnung fällt; W1 sagt bei $12$ und $16$\,Nm sofortiges Versagen vorher, weil die lineare Karte die Spannung nicht begrenzt. Die anisotrope Erweiterung mit der Hill-Vergleichs\-spannung ist als Werkzeug für die Lage des kritischen Knotens und für den Vergleich der Stufen brauchbar, nicht für die Höhe der Lebensdauer. Ihre Auslastung hängt von der Klasse der Abminderungsfaktoren ab, deren Wahl für einen Werkstoff mit richtungsabhängiger Festigkeit im Buch nicht geregelt ist. Bei W3 liegt die Vergleichsspannung zudem an bis zu 1138 Fußknoten jenseits der Isotherme. Der kritische Knoten wandert unter Hill von der Stirnseite in die Zahnmitte, wo die Faserorientierung die Querrichtung schwächt; das ist eine Aussage, die ohne die Orientierung nicht zu gewinnen ist.

Die Orientierung selbst trägt mehrere Annahmen. Sie stammt aus einem Zahn der Spritzgießsimulation, während der Index über den Umfang des Rades in Wellen von $12$ bis $15$ Zähnen um $\pm 0{,}1$ streut; welcher Zahn in das Modell übernommen wird, ist eine Eingabe. Die Simulation trifft die CT-Struktur des Berichts im Fußgrund, im Kern und im örtlichen Orientierungsgrad der Fußrundungen, nicht aber die Richtung in oberer Flanke und Kopf, wo die zugeordneten Fasern entlang der Zahnbreite laufen. Ein reales Teil mit Randschichten in der Stirnebene wäre in der Biegung steifer; mit den CT-Zonen auf demselben Netz ändert sich die Zugspannung der Fußrundung um $-8$ bis $+5$\,\%, die Drehung um $-5$ bis $-13$\,\%. Der Orientierungseffekt ist damit von der Ordnung des Körpereffekts und kleiner als der Sprung zwischen den Werkstoffstufen; die Entscheidung zwischen beiden Bildern bleibt einem Schliffbild vorbehalten.

Das FE-Modell setzt die Ritzelflanke starr und ohne Reibung an, rundet die Ritzelkopfkante mit $0{,}10$\,mm und rechnet einen Sektor aus fünf Zähnen mit einer Stellung je Analyse. Die Kopfkante bestimmt nach \secref{sec:kopfkante} den Kontaktdruck im Kantenkontakt und die Abbrüche bei scharfer Kante; die Fußspannungen an den Normpunkten ändert sie um weniger als $9$\,\%, das Globalmaximum der Zugspannung liegt bei den isotropen Stufen im verlängerten Eingriff hinter E. Die größten Auslastungen des Nachweises sitzen an den Kopfknoten unter dieser Kante und sind Kontaktzustände; sie wurden benannt und aus der Fußaussage ausgeschlossen. Die Stellungen wurden einzeln gerechnet, weil das Abwälzen mit Vorbelastung und Reibung in einem Lauf neben den Werkstoffstufen und der Kantenfrage nicht in der Zeit lag; die Lastgeschichte eines Knotens ist damit die Folge der Stellungen, nicht ein Abwälzvorgang mit Reibung.

Auf der Versuchsseite begrenzen die Belegung und die Dokumentation die Aussage. Die Bruchlinie steht auf neun Rädern mit zwei, drei und vier je Horizont; die Stufe $8$\,Nm trägt einen Lauf mit geschätzter Lastwechselzahl und einen mit erhöhtem Moment in der ersten Million Lastwechsel. Die Standardabweichung von $0{,}13$ liegt im üblichen Bereich, stützt sich aber auf eine Stichprobe, die DIN 50100 für die Streuung als unzuverlässig einstuft. Die Anrissfenster liegen nur bei $8$\,Nm beidseitig; bei $12$ und $16$\,Nm fehlt eine Kontrolle ohne Befund in der ersten Million Lastwechsel, weil die Kontrollen an die Halte des Prüfstands gebunden waren. Die Regeln für Startmoment und Kontrolle wurden im Verlauf der Reihe verschärft und gelten ab dem Lauf kst-E\_74 vollständig. Der Zähler zählt die Ritzelwelle; die Lastwechsel des Radzahns liegen um $51/52$ darunter, was unter der Auflösung der Auswertung bleibt. Die Vergleichswerkstoffe aus dem Jahr 2023 liefen auf einem anderen Prüfstand mit anderem Ritzel bei $2250$\,1/min; ihr Abstand zu den POM-Läufen des Jahres 2026 lässt sich keinem dieser Einflüsse zuordnen.

\section{Fertigungseinfluss und Fußgeometrie}
\label{sec:fertigung_fussgeometrie}

Die Prüfräder weichen vor dem Versuch von der Nenngeometrie ab, wie \secref{sec:pruefraeder_vorher} zeigt. Die Profil-Winkelabweichung ist auf beiden Flanken positiv, die Teilungs-Gesamtabweichung überschreitet die Toleranz der Stufe 7 in 40 von 58 Dateien und der Rundlauf in 55. Die Kopfkante des Rades trägt einen Bogen von $0{,}21$\,mm statt der Fase der Erzeugung. Das Rad 97345 mit den größten Teilungsabweichungen der Kerngruppe erreichte bei $8$\,Nm $25 \cdot 10^{6}$ Lastwechsel gegen geschätzt $40 \cdot 10^{6}$ des Rades 97364; ob die Teilung die Lebensdauer trägt, lässt sich mit zwei Rädern nicht entscheiden. Die versagten Zähne 31, 34 und 39 liegen außerhalb des Fensters der Zähne 36 bis 40, in dem die Teilungs-Sektorabweichung bei einzelnen Rädern am größten ist; ein Zusammenhang zwischen Teilung und Bruchzahn ist mit den vorliegenden Läufen nicht belegt. \Note{PH: Nachher-Scans der Zähne 1, 18 und 36 der Räder 97340, 97345 und 97351 gegen die Scans vor dem Versuch (M5); Schadensbilder (6.2).}

Die Fußrundungsvarianten wirken in \secref{sec:fussgeometrie} in der Reihenfolge der ebenen Studie: Die Patentkonturen senken die größte Zugspannung an C unter W1 um $25$ bis $28$\,\%, die Ellipse nach Frühe um $16$\,\% und ihr Optimum um $24$\,\%. Die Plastizität der isotropen Karte halbiert den Nutzen auf $16$ bis $21$\,\%, unter den anisotropen Karten bleiben $10$ bis $30$\,\%. Die Patentkonturen füllen den Fuß auf und machen den Zahn steifer. Dadurch steigt der Kontaktdruck an B und E um $21$ bis $30$\,\%, und der Abstand der Ritzelkopfkante zur Fußrundung schrumpft von $0{,}49$ auf $0{,}20$\,mm. Die Ellipsen legen den Fuß tiefer und verschieben den Größtwert in die Lückenmitte auf den kleinsten Krümmungsradius, wo das Netz die Zahlen begrenzt. Geometrie und Fertigung lassen sich damit getrennt beziffern. Die Fußgeometrie allein mindert die Zugspannung um bis zu $28$\,\% unter der linearen Karte. Die Fertigung allein, als Faserorientierung zwischen W1 und W3, erhöht sie um $15$ bis $20$\,\%; der Orientierungseffekt zwischen Zuordnung und CT liegt bei $-8$ bis $+5$\,\%. In der Kombination behalten die Patentkonturen unter W3 ihre Minderung von $25$ bis $33$\,\%, die Ellipsen verlieren dort an den Doppeleingriffspunkten bis zu $33$\,\%. \Note{PH: Kombination Variante mal Fertigung mit den Nachtläufen W3/W4 der Varianten und dem Stommel-Nachweis am kritischen Knoten je Variante, sobald die Kette gerechnet ist.} Für den Laufversuch bedeutet das: Eine Patentkontur verschiebt den Fuß in Richtung der Stirnseite nicht, aber sie rückt die Kopfkante des Ritzels an die Fußrundung heran; ihre Wirkung auf die Lebensdauer ist erst im Versuch an gefertigten Rädern zu belegen.

% Herkunft:
% 7.1 Absatz 1: kap5_3 (Faktoren 75/270/820: 2,671e7/3,54e5, 5,964e6/2,24e4, 1,357e6/1,65e3), kap6_3 (Anrissfenster
%   9 bis 31 % und 9 bis 17 %, obere Schranken 55 bis 93 %), kap3_8 (N_A ≤ N_K ≤ N_B), ADR-118 Amendment 2026-10-18
%   Abend Punkt 1 (dreistufig).
% 7.1 Absatz 2: Steigungen aus kap5_5 (σ1 an C) und woehler.md: W2 lg(1,357e6/2,671e7) / lg(38,27/48,90) =
%   −1,294/−0,1065 = 12,2; W4 −1,294/lg(48,58/65,23) = −1,294/−0,128 = 10,1; W1 −1,294/lg(56,09/82,87) =
%   −1,294/−0,1695 = 7,6; W3 −1,294/lg(65,07/97,35) = −1,294/−0,175 = 7,4; synthetische Linie: stommel_W2_QS
%   bericht (σ_W 20,0 MPa, k 14,77 bis 15,77, Tabelle 6.3 15,65 bis 20,44); E8 offen (Roadmap).
% 7.1 Absatz 3: ADR-118 (Isotropievorschrift, 0°-Kurve als Referenz; kritischer Knoten Druckseite R ≈ −0,4,
%   z −6,0 mm); kap4_5 (84,1/87,8 °C bei 16 Nm); 50 Hz = 3000 1/min; kap4_3 (Raten); kap5_3 (S-Life Faktor 2,45
%   bis 3,11, Kappe); Lesart offen (E4).
% 7.1 Absatz 4: kap5_5 Platzhalter; kap6_3 (Netze 0,60 bis 0,74); kap6_4 (1,6 bis 2,1; PEEK Flanke, MA-0296 S. 66).
% 7.2 Absatz 1: kap5_3 (W2 besteht statisch, W1 sofort bei 12/16 Nm; Hill: Klasse, 1138 Knoten jenseits der
%   Isotherme, kritischer Knoten in Zahnmitte); ADR-118 E7 (Klasse nicht geregelt, Buch S. 346).
% 7.2 Absatz 2: kap5_1 (Wellen 12 bis 15 Zähne, ±0,1; CT gegen Zuordnung −8 bis +5 % σ1, −5 bis −13 % Drehung;
%   Schliffbild offen); ADR-119 Consequences.
% 7.2 Absatz 3: kap3_6 (starre Ritzelfläche, keine Reibung, Kante 0,10 mm, Sektor fünf Zähne, eine Stellung je
%   Analyse); kap4_2 (Kante, Abbrüche); kap5_2 (Fußspannung an B/E −5 bis −9 %, Maximum hinter E); kap5_3 (Kopfknoten);
%   Plan 2026-10-14 Kapitel 8 (Abwälzen als Ausblick).
% 7.2 Absatz 4: kap6_1 und kap6_3 (Belegung 2/3/4, geschätzt, erhöhtes Moment, s 0,129, DIN 50100 S. 30 n < 10,
%   Anrissfenster, Regeln ab kst-E_74, Zähler Ritzelwelle, 2023 anderer Prüfstand/Ritzel/Drehzahl); kap6_4.
% 7.3 Absatz 1: kap4_4 (f_Hα positiv, F_p 40 von 58, F_r 55, Bogen 0,21 mm, 97345 F_p 136/123 µm, Fenster 36 bis
%   40, F_pk); kap6_1 (97345 24,96e6, 97364 40e6 geschätzt; Zähne 31, 34, 39, 40, 41); M5 offen.
% 7.3 Absatz 2: kap5_4 (Tabelle C: Ellipse −15,9/−24,4 %, Bézier −25,2/−27,6 %, bionisch −27,4 % unter W1; W2 −6,0
%   bis −20,5 %; W3 −4,3 bis −29,1 %; Patentkonturen A bis E W1 −25 bis −35 %, W3 −25 bis −33 %; Kontaktdruck B/E
%   +21 bis +30 %; Abstand Kopfkante 0,49 → 0,20 mm; Ellipsen an A/B/D/E bis +33 % unter W3, Netzgrenze FE-25);
%   Fertigung allein: kap5_2 (W1 → W3 +15 bis +20 %) und kap5_1 (−8 bis +5 %); Nutzerwunsch 2026-10-18 (einzeln je
%   Seite und in Kombination); Nachtläufe W3/W4 der Varianten offen (Serverraum Montag).
% Bildwunsch: keiner.
% Bibwunsch: keine neuen Schlüssel.
% Nomenklatur: N_A, N_K, N_B, k, R, sigma_1 vorhanden.
% Offen: E8-Zahlen; VDI/RIKOR/Workbench; Lesart der Kappe; Nachher-Scans; Kombination mit den Nachtläufen;
%   Schadensbilder.
